\chapter{Introduzione}
\label{cap:sdd-introduzione}

\section{Scopo del documento}

Questo documento descrive come Tech4All è costruito. Muove dai requisiti non
funzionali del RAD per ricavarne gli obiettivi di progettazione, decompone il
sistema in sottosistemi, ne definisce l'interfaccia, e scende fino alla
progettazione delle classi che li realizzano.

Rispetto all'impostazione di Bruegge e Dutoit~\cite{bruegge2010} il documento
accorpa il System Design Document e l'Object Design Document. La scelta è
motivata dalla scala del sistema: un ODD autonomo, per un progetto di queste
dimensioni, sarebbe in gran parte frontespizio, glossario e ripetizione delle
premesse. L'object design occupa quindi il \cref{cap:sdd-object-design}, con
tutti i contenuti che il testo di riferimento gli attribuisce.

\section{Obiettivi di progettazione}
\label{sec:sdd-design-goals}

Gli obiettivi di progettazione (\term{design goal}) traducono i requisiti non
funzionali in criteri che guidano le decisioni tecniche. Ciascuno è derivato
da uno o più requisiti del RAD; la corrispondenza completa è
nell'appendice~\ref{cap:sdd-app-tracciabilita}.

\begin{longtable}{@{}l L{4.4cm} L{7.4cm}@{}}
\toprule
\textbf{Id} & \textbf{Obiettivo} & \textbf{Significato operativo} \\
\midrule
\endfirsthead
\toprule
\textbf{Id} & \textbf{Obiettivo} & \textbf{Significato operativo} \\
\midrule
\endhead
\bottomrule
\endlastfoot

\id{DG\_1} & Fiducia nel confine &
Ciò che arriva dall'esterno non è mai considerato attendibile: né i dati,
né l'identità dichiarata, né i nomi dei file. Il sistema decide sulla base
di ciò che possiede, non di ciò che gli viene detto. \\

\id{DG\_2} & Un solo luogo per ogni decisione &
Ogni regola — un vincolo di dominio, una soglia, un permesso — è definita in
un punto solo. Regole duplicate divergono. \\

\id{DG\_3} & Verificabilità &
Ogni componente deve poter essere esercitato in isolamento, senza database né
rete. Ciò che non si riesce a provare facilmente è, di norma, mal progettato. \\

\id{DG\_4} & Uniformità dei sottosistemi &
I sottosistemi hanno la stessa struttura interna e le stesse convenzioni.
Chi ne conosce uno sa orientarsi negli altri. \\

\id{DG\_5} & Coerenza dei dati &
Le operazioni che toccano più entità sono atomiche; gli invarianti che il
database può garantire sono affidati al database. \\

\id{DG\_6} & Costo di esercizio contenuto &
Nessuna dipendenza da servizi a pagamento, nessuna infrastruttura oltre un
processo applicativo e un DBMS. \\

\end{longtable}

\section{Compromessi}
\label{sec:sdd-tradeoff}

Gli obiettivi di progettazione entrano in conflitto fra loro. I compromessi
che seguono sono le decisioni prese quando non era possibile soddisfarli
tutti, con la ragione della scelta.

\begin{longtable}{@{}L{3.6cm} L{4.2cm} L{6.4cm}@{}}
\toprule
\textbf{Compromesso} & \textbf{Alternativa scartata} & \textbf{Motivazione} \\
\midrule
\endfirsthead
\toprule
\textbf{Compromesso} & \textbf{Alternativa scartata} & \textbf{Motivazione} \\
\midrule
\endhead
\bottomrule
\endlastfoot

Sicurezza contro immediatezza della risposta &
Correggere il quiz nel browser, evitando un giro di rete &
La correzione locale renderebbe le soluzioni visibili prima della consegna,
annullando il valore della verifica (\id{OB\_6}). Il costo è una richiesta
HTTP: trascurabile rispetto a ciò che si guadagna. \\

Uniformità contro concisione &
Servizi ridotti a semplici passacarte verso i DAO &
Alcuni servizi si limitano a delegare, e potrebbero essere eliminati. Sono
stati mantenuti perché la loro assenza renderebbe i sottosistemi diversi fra
loro senza un motivo di dominio (\id{DG\_4}). \\

Coerenza contro indipendenza dal DBMS &
Calcolare la valutazione media nel codice applicativo &
La media è mantenuta da trigger. Il valore resta corretto anche di fronte a
scritture concorrenti, al prezzo di una logica che vive nel database e non è
verificata dai test unitari (\id{DG\_5}). \\

Semplicità contro completezza funzionale &
Recupero della password via email &
Realizzarlo richiederebbe un servizio di posta transazionale e un dominio
verificato, contro \id{DG\_6}. Il sistema rinuncia alla funzionalità e lo
dichiara apertamente. \\

Verificabilità contro fedeltà dei test &
Test di integrazione su un database reale &
La suite sostituisce il livello di persistenza con dei doppi: gira ovunque in
pochi secondi (\id{DG\_3}), ma non verifica le query SQL, che restano coperte
solo dalle prove manuali di sistema. \\

\end{longtable}

\begin{nota}
L'ultimo compromesso è quello con il rischio residuo più alto ed è
esplicitamente dichiarato anche nel Test Document: un errore in una query
verrebbe intercettato solo all'esecuzione. È stato accettato perché il
livello di persistenza è sottile e privo di logica condizionale, mentre il
costo di un database nella catena di integrazione continua sarebbe permanente.
\end{nota}

\section{Definizioni e riferimenti}

I termini del dominio e quelli tecnici sono definiti nel
\cref{cap:glossario}, condiviso con gli altri documenti del progetto. Il
documento presuppone la lettura del RAD, di cui riprende identificativi di
requisiti e casi d'uso senza ridefinirli.
